%------------------------------------------------------------------------------%
\documentclass[a4paper,12pt,fleqn]{article}

\usepackage{calculo_varias_variaveis-1}

\ShowAnswers

%------------------------------------------------------------------------------%
\begin{document}

\makeHeader{CFVVI}{Prova 3 -- Turma 4}{14/07/2025}

% 01
%------------------------------------------------------------------------------%
\exercise{20}
Converta a equação para coordenadas cartesianas e mostre que a solução é uma cônica
\[
  r = \frac{12}{1 - 3\cos\theta}
\]
\clearpagequestiononly

\begin{answer}
  Queremos mostrar que a equação, escrita em coordenadas cartesianas, $(x, y)$,
  é uma expressão com termos de segundo grau em $x$ e $y$
  \begin{align*}
    r                       & = \frac{12}{1 - 3\cos\theta} \\
    r(1 - 3\cos\theta)      & = 12                         \\
    r - 3r\cos\theta        & = 12                         \\
    \sqrt{x^2 + y^2} - 3x   & = 12                         \\
    \sqrt{x^2 + y^2}        & = 12 + 3x                    \\
    x^2 + y^2               & = 144 + 72x + 9x^2           \\
    -8x^2 + y^2 - 72x - 144 & = 0                          \\
    8x^2 - y^2 + 72x + 144  & = 0
  \end{align*}
\end{answer}

% 02
%------------------------------------------------------------------------------%
\exercise{20}
Encontre as raízes cúbicas de \( z = 27\left(\frac{1}{2} - i\frac{\sqrt{3}}{2}\right) \)
\clearpagequestiononly

\begin{answer}
    Converter para forma polar
    \begin{align*}
        x       & = \Re(z) = \frac{27}{2}         \\
        y       & = \Im(z) = \frac{-27\sqrt{3}}{2} \\
        |z|     & = \sqrt{x^2+y^2}
                  = \sqrt{\left(\frac{27}{2}\right)^2 + \left(\frac{-27\sqrt{3}}{2}\right)^2}
                  = \sqrt{\frac{27^2}{2^2} + \frac{27^2\times 3}{2^2}}
                  = \sqrt{\frac{27^2\times 4}{2^2}}
                  = \frac{27\times 2}{2}
                  = 27         \\
        \varphi & = \arg(z)
                  = \arctan\left(\frac{y}{x}\right)
                  = \arctan\left(\frac{\dfrac{-27\sqrt{3}}{2}}{\dfrac{27}{2}}\right)
                  = \arctan\left(-\sqrt{3}\right)
                  = - \frac{\pi}{3} \\
        z       & = 27 \cis{\frac{-\pi}{3}}
        = 27 e^{-i\pi/3}
      \end{align*}

    Aplicar a fórmula de De Moivre para as raízes cúbicas
    \[
      u_k
      = \sqrt[3]{27}\cis{\frac{-\pi/3 + 2k\pi}{3}}
      \qquad\qquad k = 0,1,2
    \]
    Calculando os argumentos das raízes
    \begin{align*}
        \varphi_0
        & = \frac{-\pi/3 + 2\times 0\pi}{3}
          = \frac{-\pi}{9} \\
        \varphi_1
        & = \frac{-\pi/3 + 2\times 1\pi}{3}
          = \frac{-\pi/3 + 6\pi/3}{3}
          = \frac{5\pi}{9} \\
        \varphi_2
        & = \frac{-\pi/3 + 2\times 2\pi}{3}
          = \frac{-\pi/3 + 12\pi/3}{3}
          = \frac{11\pi}{9}
      \end{align*}
    As raízes são
    \begin{align*}
        u_0 & = 3\cis{\frac{ -\pi}{9}} = 3e^{\sfrac{- i\pi}{9}} \\
        u_1 & = 3\cis{\frac{ 5\pi}{9}} = 3e^{\sfrac{5 i\pi}{9}} \\
        u_2 & = 3\cis{\frac{11\pi}{9}} = 3r^{\sfrac{11i\pi}{9}}
      \end{align*}
\end{answer}

% 03
%------------------------------------------------------------------------------%
\exercise{30}
Encontre e classifique os pontos críticos da função
\(
  f(x, y) = y^3 + 3x^2y - 6x^2 - 6y^2 + 2
\)
\clearpagequestiononly

\begin{answer}
\setlength\columnsep{3em}
\begin{multicols*}{2}
\raggedcolumns
  Calculando as \textbf{derivadas primeiras} de $f$
  \begin{align*}
    \D{f}{x}
    & = \D{}{x}\left(y^3 + 3x^2y - 6x^2 - 6y^2 + 2\right) \\
    & = 0 + 3\times 2xy - 6\times 2x - 0 + 0 \\
    & = 6xy - 12x \\
    & = 6x(y - 2)
  \end{align*}
  \begin{align*}
    \D{f}{y}
    & = \D{}{y}\left(y^3 + 3x^2y - 6x^2 - 6y^2 + 2\right) \\
    & = 3y^2 + 3x^2 - 0 - 6\times 2y + 0 \\
    & = 3y^2 + 3x^2 - 12y \\
    & = 3(y^2 - 4y + x^2)
  \end{align*}

  Encontrando os \textbf{pontos críticos}.
  As derivadas parciais de $f$ existem em todo o plano, temos então que resolver o sistema
  \begin{align*}
    \D{f}{x} & = 6x(y - 2) = 0 \\
    \D{f}{y} & = 3(y^2 - 4y + x^2) = 0
  \end{align*}

  Da primeira equação temos que $x=0$ ou $y=2$

  Se $x=0$ temos
  \begin{align*}
    3(y^2 - 4y + x^2) & = 0 \\
    y(y - 4)          & = 0
  \end{align*}
  cujas soluções são $y=0$ ou $y=4$. Temos então os pontos críticos
  \[
    P_1 = (0, 0)
    \quad\text{e}\quad
    P_2 = (0, 4)
  \]

  Se $y=2$ temos
  \begin{align*}
    3(y^2 - 4y + x^2)     & = 0     \\
    2^2 - 4\times 2 + x^2 & = 0     \\
    4 - 8 + x^2           & = 0     \\
    x^2                   & = 4     \\
    x                     & = \pm 2
  \end{align*}
  cujas soluções são $x=2$ ou $x=-2$. Temos então os pontos críticos
  \[
    P_3 = (2, 2)
    \quad\text{e}\quad
    P_4 = (-2, 2)
  \]

  Para \textbf{classificar} os pontos críticos precisamos
  do determinante da Hessiana
  \begin{align*}
    \DS{f}{x}    & = \D{}{x}\left(6xy - 12x\right)         = 6y - 12 \\
    \DS{f}{y}    & = \D{}{y}\left(3y^2 + 3x^2 - 12y\right) = 6y - 12 \\
    \DM{f}{y}{x} & = \D{}{y}\left(6xy - 12x\right)         = 6x
  \end{align*}
  \begin{align*}
    D(x, y)
    & = f_{xx}(x, y) f_{yy}(x, y) - \left(f_{xy}(x, y)\right)^2 \\
    & = (6y-12)(6y-12) - (6x)^2 \\
    & = (6y-12)^2 - 36x^2
  \end{align*}

  Classificando o ponto $P_1$
  \begin{align*}
    D\left(0, 0\right)
    & = (6\times 0-12)^2 - 36\times 0^2 \\
    & = (-12)^2 > 0 \\
    f_{xx}\left(0, 0\right)
    & = 6\times 0 - 12 = -12 < 0
  \end{align*}
  O ponto $P_1$ é um máximo local

  Classificando o ponto $P_2$
  \begin{align*}
    D\left(0, 4\right)
    & = (6\times 4-12)^2 - 36\times 0^2 \\
    & = (24-12)^2 > 0 \\
    f_{xx}\left(0, 4\right)
    & = 6\times 4 - 12 = 12 > 0
  \end{align*}
  O ponto $P_2$ é um mínimo local

  Classificando o ponto $P_3$
  \begin{align*}
    D\left(2, 2\right)
    & = (6\times 2-12)^2 - 36\times 2^2 \\
    & = (12-12)^2 - 36\times 4 \\
    & = -36\times 4 < 0
  \end{align*}
  $P_3$ é um ponto de sela

  Classificando o ponto $P_4$
  \begin{align*}
    D\left(-2, 2\right)
    & = (6\times 2-12)^2 - 36\times(-2)^2 \\
    & = (12-12)^2 - 36\times 4 \\
    & = -36\times 4 < 0
  \end{align*}
  $P_4$ é um ponto de sela
\end{multicols*}
\end{answer}

% 04
%------------------------------------------------------------------------------%
\exercise{30}
Encontre o valor mínimo de
\(
  e^{xy}
\)
sujeitos a
\(
  x^3 + y^3 = 16
\)
\clearpagequestiononly

\begin{answer}
  \setlength\columnsep{3em}
  \begin{multicols}{2}
    Queremos encontrar o valor mínimo de $e^{xy}$
    mas como a função exponencial é crescente podemos localizar  ponto
    onde o mínimo acontece usando a função
    \[
      f(x,y) = xy
    \]
    sujeitos a restrição
    \[
      g(x, y) = x^3 + y^3 = 16
    \]
    Gradiente de $f$
    \[
      \D{f}{x}
      = \D{}{x}\left(xy\right)
      = y
    \]
    \[
      \D{f}{y}
      = \D{}{y}\left(xy\right)
      = x
    \]
    \[
      \nabla f(x, y) = \vetor{y}{x}
    \]
    Gradiente de $g$
    \[
      \D{g}{x}
      = \D{}{x}\left(x^3 + y^3\right)
      = 3x^2
    \]
    \[
      \D{g}{y}
      = \D{}{y}\left(x^3 + y^3\right)
      = 3y^2
    \]
    \[
      \nabla g(x, y) = \vetor{3x^2}{3y^2}
    \]
    Precisamos resolver o sistema
    \begin{align*}
      y         & = \lambda 3x^2 \\
      x         & = \lambda 3y^2 \\
      x^3 + y^3 & = 16
    \end{align*}
    Que podemos manipular as duas primeiras equações para obter
    \begin{align*}
      \dfrac{xy}{3} & = \lambda x^3 \\
      \dfrac{xy}{3} & = \lambda y^3
    \end{align*}
    Portanto,
    \(
      \lambda x^3 = \lambda y^3
    \) \\
    Assim $\lambda=0$ ou $x=y$ \\
    Se $\lambda=0$ as duas primeiras equações
    determinam $x=0$ e $y=0$ o que é incompatível com a
    terceira equação. Portanto, sabemos que $\lambda\neq 0$.

    Substituindo $x=y$ na terceira equação
    \begin{align*}
      x^3 + y^3 & = 16 \\
      2x^3      & = 16 \\
      x^3       & = 8  \\
      x         & = 2
    \end{align*}
    Portanto o único ponto candidato a mínimo é $(x, y) = (2, 2)$.

    Avaliando $e^{xy}$ no ponto temos $\min = e^4$
  \end{multicols}
\end{answer}

%------------------------------------------------------------------------------%
\end{document}
%------------------------------------------------------------------------------%
